\section{Open immersions and (pre)sheaves}
\label{section-open-immersions}

\noindent
Let $X$ be a topological space.
Let $j : U \to X$ be the inclusion of an open subset $U$ into $X$.
In Section \ref{section-presheaves-functorial} we have defined
functors $j_*$ and $j^{-1}$ such that $j_*$ is right adjoint to
$j^{-1}$. It turns out that for an open immersion there is a left adjoint
for $j^{-1}$, which we will denote $j_!$. First we point out that
$j^{-1}$ has a particularly simple description in the case of
an open immersion.

\begin{lemma}
\label{lemma-j-pullback}
Let $X$ be a topological space.
Let $j : U \to X$ be the inclusion of an open subset $U$ into $X$.
\begin{enumerate}
\item Let $\mathcal{G}$ be a presheaf of sets on $X$.
The presheaf $j_p\mathcal{G}$
(see Section \ref{section-presheaves-functorial}) is given by the rule
$V \mapsto \mathcal{G}(V)$ for $V \subset U$ open.
\item Let $\mathcal{G}$ be a sheaf of sets on $X$.
The sheaf $j^{-1}\mathcal{G}$ is given by the rule
$V \mapsto \mathcal{G}(V)$ for $V \subset U$ open.
\item For any point $u \in U$ and any sheaf $\mathcal{G}$ on $X$
we have a canonical identification of stalks
$$
j^{-1}\mathcal{G}_u = (\mathcal{G}|_U)_u = \mathcal{G}_u.
$$
\item On the category of presheaves of $U$ we have $j_pj_* = \text{id}$.
\item On the category of sheaves of $U$ we have $j^{-1}j_* = \text{id}$.
\end{enumerate}
The same description holds for (pre)sheaves of abelian groups,
(pre)sheaves of algebraic structures, and (pre)sheaves of modules.
\end{lemma}

\begin{proof}
The colimit in the definition of $j_p\mathcal{G}(V)$
is over collection of all $W \subset X$ open such that $V \subset W$
ordered by reverse inclusion.
Hence this has a largest element, namely $V$. This proves (1).
And (2) follows because the assignment $V \mapsto \mathcal{G}(V)$
for $V \subset U$ open is clearly a sheaf if $\mathcal{G}$ is a
sheaf. Assertion (3) follows from (2) since the collection
of open neighbourhoods of $u$ which are contained in $U$ is cofinal
in the collection of all open neighbourhoods of $u$ in $X$.
Parts (4) and (5) follow by computing
$j^{-1}j_*\mathcal{F}(V) = j_*\mathcal{F}(V) = \mathcal{F}(V)$.

\medskip\noindent
The exact same arguments work for (pre)sheaves of abelian groups
and (pre)sheaves of algebraic structures.
\end{proof}

\begin{definition}
\label{definition-restriction}
Let $X$ be a topological space.
Let $j : U \to X$ be the inclusion of an open subset.
\begin{enumerate}
\item Let $\mathcal{G}$ be a presheaf of sets, abelian groups or
algebraic structures on $X$. The presheaf $j_p\mathcal{G}$ described
in Lemma \ref{lemma-j-pullback} is called
the {\it restriction of $\mathcal{G}$ to $U$} and denoted $\mathcal{G}|_U$.
\item Let $\mathcal{G}$ be a sheaf of sets on $X$, abelian groups or
algebraic structures on $X$. The sheaf $j^{-1}\mathcal{G}$ is called
the {\it restriction of $\mathcal{G}$ to $U$} and denoted $\mathcal{G}|_U$.
\item If $(X, \mathcal{O})$ is a ringed space, then the pair
$(U, \mathcal{O}|_U)$ is called the
{\it open subspace of $(X, \mathcal{O})$ associated to $U$}.
\item If $\mathcal{G}$ is a presheaf of $\mathcal{O}$-modules
then $\mathcal{G}|_U$ together with the multiplication map
$\mathcal{O}|_U \times \mathcal{G}|_U \to \mathcal{G}|_U$
(see Lemma \ref{lemma-pullback-module})
is called the {\it restriction of $\mathcal{G}$ to $U$}.
\end{enumerate}
\end{definition}

\noindent
We leave a definition of the restriction of presheaves
of modules to the reader. Ok, so in this section we will
discuss a left adjoint to the restriction functor.
Here is the definition in the case of (pre)sheaves
of sets.

\begin{definition}
\label{definition-j-shriek}
Let $X$ be a topological space.
Let $j : U \to X$ be the inclusion of an open subset.
\begin{enumerate}
\item Let $\mathcal{F}$ be a presheaf of sets on $U$. We define
the {\it extension of $\mathcal{F}$ by the empty set $j_{p!}\mathcal{F}$}
to be the presheaf of sets on $X$ defined by the rule
$$
j_{p!}\mathcal{F}(V) =
\left\{
\begin{matrix}
\emptyset & \text{if} & V \not \subset U \\
\mathcal{F}(V) & \text{if} & V \subset U
\end{matrix}
\right.
$$
with obvious restriction mappings.
\item Let $\mathcal{F}$ be a sheaf of sets on $U$. We define
the {\it extension of $\mathcal{F}$ by the empty set $j_!\mathcal{F}$}
to be the sheafification of the presheaf $j_{p!}\mathcal{F}$.
\end{enumerate}
\end{definition}

\begin{lemma}
\label{lemma-j-shriek}
Let $X$ be a topological space.
Let $j : U \to X$ be the inclusion of an open subset.
\begin{enumerate}
\item The functor $j_{p!}$ is a left adjoint to the
restriction functor $j_p$ (see Lemma \ref{lemma-j-pullback}).
\item The functor $j_!$ is a left adjoint to restriction,
in a formula
$$
\Mor_{\Sh(X)}(j_!\mathcal{F}, \mathcal{G})
=
\Mor_{\Sh(U)}(\mathcal{F}, j^{-1}\mathcal{G})
=
\Mor_{\Sh(U)}(\mathcal{F}, \mathcal{G}|_U)
$$
bifunctorially in $\mathcal{F}$ and $\mathcal{G}$.
\item Let $\mathcal{F}$ be a sheaf of sets on $U$.
The stalks of the sheaf $j_!\mathcal{F}$ are described
as follows
$$
j_{!}\mathcal{F}_x =
\left\{
\begin{matrix}
\emptyset & \text{if} & x \not \in U \\
\mathcal{F}_x & \text{if} & x \in U
\end{matrix}
\right.
$$
\item On the category of presheaves of $U$ we have $j_pj_{p!} = \text{id}$.
\item On the category of sheaves of $U$ we have $j^{-1}j_! = \text{id}$.
\end{enumerate}
\end{lemma}

\begin{proof}
To map $j_{p!}\mathcal{F}$ into $\mathcal{G}$
it is enough to map $\mathcal{F}(V) \to \mathcal{G}(V)$
whenever $V \subset U$ compatibly with restriction
mappings. And by Lemma \ref{lemma-j-pullback}
the same description holds for maps
$\mathcal{F} \to \mathcal{G}|_U$.
The adjointness of $j_!$ and restriction follows
from this and the properties of sheafification.
The identification of stalks is obvious from the
definition of the extension by the empty set
and the definition of a stalk.
Statements (4) and (5) follow by computing the
value of the sheaf on any open of $U$.
\end{proof}

\noindent
Note that if $\mathcal{F}$ is a sheaf
of abelian groups on $U$, then in general $j_!\mathcal{F}$ as
defined above, is not a sheaf of abelian groups, for example
because some of its stalks are empty (hence not abelian groups
for sure). Thus we need to modify the definition of
$j_!$ depending on the type of sheaves we consider.
The reason for choosing the empty set in the definition of the
extension by the empty set, is that it is the initial object
in the category of sets. Thus in the case of abelian groups
we use $0$ (and more generally for sheaves with values in
any abelian category).

\begin{definition}
\label{definition-j-shriek-structures}
Let $X$ be a topological space.
Let $j : U \to X$ be the inclusion of an open subset.
\begin{enumerate}
\item Let $\mathcal{F}$ be an abelian presheaf on $U$.
We define the {\it extension $j_{p!}\mathcal{F}$ of $\mathcal{F}$ by $0$}
to be the abelian presheaf on $X$ defined by the rule
$$
j_{p!}\mathcal{F}(V) =
\left\{
\begin{matrix}
0 & \text{if} & V \not \subset U \\
\mathcal{F}(V) & \text{if} & V \subset U
\end{matrix}
\right.
$$
with obvious restriction mappings.
\item Let $\mathcal{F}$ be an abelian sheaf on $U$. We define
the {\it extension $j_!\mathcal{F}$ of $\mathcal{F}$ by $0$}
to be the sheafification of the abelian presheaf $j_{p!}\mathcal{F}$.
\item Let $\mathcal{C}$ be a category having an initial object $e$.
Let $\mathcal{F}$ be a presheaf on $U$ with values in $\mathcal{C}$.
We define the {\it extension $j_{p!}\mathcal{F}$ of $\mathcal{F}$ by $e$}
to be the presheaf on $X$ with values in $\mathcal{C}$ defined by the
rule
$$
j_{p!}\mathcal{F}(V) =
\left\{
\begin{matrix}
e & \text{if} & V \not \subset U \\
\mathcal{F}(V) & \text{if} & V \subset U
\end{matrix}
\right.
$$
with obvious restriction mappings.
\item Let $(\mathcal{C}, F)$ be a type of algebraic structure
such that $\mathcal{C}$ has an initial object $e$.
Let $\mathcal{F}$ be a sheaf of algebraic structures on $U$
(of the give type). We define the
{\it extension $j_!\mathcal{F}$ of $\mathcal{F}$ by $e$}
to be the sheafification of the presheaf $j_{p!}\mathcal{F}$
defined above.
\item Let $\mathcal{O}$ be a presheaf of rings on $X$.
Let $\mathcal{F}$ be a presheaf of $\mathcal{O}|_U$-modules.
In this case we define the {\it extension by $0$}
to be the presheaf of $\mathcal{O}$-modules which is equal to
$j_{p!}\mathcal{F}$ as an abelian presheaf endowed with
the multiplication map
$\mathcal{O} \times j_{p!}\mathcal{F} \to j_{p!}\mathcal{F}$.
\item Let $\mathcal{O}$ be a sheaf of rings on $X$.
Let $\mathcal{F}$ be a sheaf of $\mathcal{O}|_U$-modules.
In this case we define the {\it extension by $0$}
to be the $\mathcal{O}$-module which is equal to
$j_!\mathcal{F}$ as an abelian sheaf endowed with
the multiplication map $\mathcal{O} \times j_!\mathcal{F} \to j_!\mathcal{F}$.
\end{enumerate}
\end{definition}

\noindent
It is true that one can define $j_!$ in the setting of sheaves
of algebraic structures (see below). However, it depends on the
type of algebraic structures involved what the resulting
object is. For example, if $\mathcal{O}$ is a sheaf of rings
on $U$, then $j_{!, rings}\mathcal{O} \not = j_{!, abelian}\mathcal{O}$
since the initial object in the category of rings
is $\mathbf{Z}$ and the initial object in the category
of abelian groups is $0$. In particular the functor $j_!$
{\it does not commute with taking underlying sheaves of sets},
in contrast to what we have seen so far! We separate out the case
of (pre)sheaves of abelian groups, (pre)sheaves of algebraic structures
and (pre)sheaves of modules as usual.

\begin{lemma}
\label{lemma-j-shriek-abelian}
Let $X$ be a topological space.
Let $j : U \to X$ be the inclusion of an open subset.
Consider the functors of restriction and extension
by $0$ for abelian (pre)sheaves.
\begin{enumerate}
\item The functor $j_{p!}$ is a left adjoint to the
restriction functor $j_p$ (see Lemma \ref{lemma-j-pullback}).
\item The functor $j_!$ is a left adjoint to restriction,
in a formula
$$
\Mor_{\textit{Ab}(X)}(j_!\mathcal{F}, \mathcal{G})
=
\Mor_{\textit{Ab}(U)}(\mathcal{F}, j^{-1}\mathcal{G})
=
\Mor_{\textit{Ab}(U)}(\mathcal{F}, \mathcal{G}|_U)
$$
bifunctorially in $\mathcal{F}$ and $\mathcal{G}$.
\item Let $\mathcal{F}$ be an abelian sheaf on $U$.
The stalks of the sheaf $j_!\mathcal{F}$ are described
as follows
$$
j_{!}\mathcal{F}_x =
\left\{
\begin{matrix}
0 & \text{if} & x \not \in U \\
\mathcal{F}_x & \text{if} & x \in U
\end{matrix}
\right.
$$
\item On the category of abelian presheaves of $U$
we have $j_pj_{p!} = \text{id}$.
\item On the category of abelian sheaves of $U$
we have $j^{-1}j_! = \text{id}$.
\end{enumerate}
\end{lemma}

\begin{proof}
Omitted.
\end{proof}

\begin{lemma}
\label{lemma-j-shriek-structures}
Let $X$ be a topological space.
Let $j : U \to X$ be the inclusion of an open subset.
Let $(\mathcal{C}, F)$ be a type of algebraic structure
such that $\mathcal{C}$ has an initial object $e$.
Consider the functors of restriction and extension
by $e$ for (pre)sheaves of algebraic structure defined above.
\begin{enumerate}
\item The functor $j_{p!}$ is a left adjoint to the
restriction functor $j_p$ (see Lemma \ref{lemma-j-pullback}).
\item The functor $j_!$ is a left adjoint to restriction,
in a formula
$$
\Mor_{\Sh(X, \mathcal{C})}(j_!\mathcal{F}, \mathcal{G})
=
\Mor_{\Sh(U, \mathcal{C})}(\mathcal{F}, j^{-1}\mathcal{G})
=
\Mor_{\Sh(U, \mathcal{C})}(\mathcal{F}, \mathcal{G}|_U)
$$
bifunctorially in $\mathcal{F}$ and $\mathcal{G}$.
\item Let $\mathcal{F}$ be a sheaf on $U$.
The stalks of the sheaf $j_!\mathcal{F}$ are described
as follows
$$
j_{!}\mathcal{F}_x =
\left\{
\begin{matrix}
e & \text{if} & x \not \in U \\
\mathcal{F}_x & \text{if} & x \in U
\end{matrix}
\right.
$$
\item On the category of presheaves of algebraic structures on $U$
we have $j_pj_{p!} = \text{id}$.
\item On the category of sheaves of algebraic structures on $U$
we have $j^{-1}j_! = \text{id}$.
\end{enumerate}
\end{lemma}

\begin{proof}
Omitted.
\end{proof}

\begin{lemma}
\label{lemma-j-shriek-modules}
Let $(X, \mathcal{O})$ be a ringed space.
Let $j : (U, \mathcal{O}|_U) \to (X, \mathcal{O})$
be an open subspace.
Consider the functors of restriction and extension
by $0$ for (pre)sheaves of modules defined above.
\begin{enumerate}
\item The functor $j_{p!}$ is a left adjoint to restriction,
in a formula
$$
\Mor_{\textit{PMod}(\mathcal{O})}(j_{p!}\mathcal{F}, \mathcal{G})
=
\Mor_{\textit{PMod}(\mathcal{O}|_U)}(\mathcal{F}, \mathcal{G}|_U)
$$
bifunctorially in $\mathcal{F}$ and $\mathcal{G}$.
\item The functor $j_!$ is a left adjoint to restriction,
in a formula
$$
\Mor_{\textit{Mod}(\mathcal{O})}(j_!\mathcal{F}, \mathcal{G})
=
\Mor_{\textit{Mod}(\mathcal{O}|_U)}(\mathcal{F}, \mathcal{G}|_U)
$$
bifunctorially in $\mathcal{F}$ and $\mathcal{G}$.
\item Let $\mathcal{F}$ be a sheaf of $\mathcal{O}$-modules on $U$.
The stalks of the sheaf $j_!\mathcal{F}$ are described
as follows
$$
j_{!}\mathcal{F}_x =
\left\{
\begin{matrix}
0 & \text{if} & x \not \in U \\
\mathcal{F}_x & \text{if} & x \in U
\end{matrix}
\right.
$$
\item On the category of sheaves of $\mathcal{O}|_U$-modules on $U$
we have $j^{-1}j_! = \text{id}$.
\end{enumerate}
\end{lemma}

\begin{proof}
Omitted.
\end{proof}

\noindent
Note that by the lemmas above, both the functors
$j_*$ and $j_!$ are fully faithful embeddings of
the category of sheaves on $U$ into the category
of sheaves on $X$. It is only true for the functor
$j_!$ that one can easily describe the essential
image of this functor.

\begin{lemma}
\label{lemma-equivalence-categories-open}
Let $X$ be a topological space.
Let $j : U \to X$ be the inclusion of an open subset.
The functor
$$
j_! : \Sh(U) \longrightarrow \Sh(X)
$$
is fully faithful. Its essential image consists exactly
of those sheaves $\mathcal{G}$ such that
$\mathcal{G}_x = \emptyset$ for all $x \in X \setminus U$.
\end{lemma}

\begin{proof}
Fully faithfulness follows formally from $j^{-1} j_! = \text{id}$.
We have seen that any sheaf in the image of the functor has
the property on the stalks mentioned in the lemma. Conversely, suppose
that $\mathcal{G}$ has the indicated property.
Then it is easy to check that
$$
j_! j^{-1} \mathcal{G} \to \mathcal{G}
$$
is an isomorphism on all stalks and hence an isomorphism.
\end{proof}

\begin{lemma}
\label{lemma-equivalence-categories-open-abelian}
Let $X$ be a topological space.
Let $j : U \to X$ be the inclusion of an open subset.
The functor
$$
j_! : \textit{Ab}(U) \longrightarrow \textit{Ab}(X)
$$
is fully faithful. Its essential image consists exactly
of those sheaves $\mathcal{G}$ such that
$\mathcal{G}_x = 0$ for all $x \in X \setminus U$.
\end{lemma}

\begin{proof}
Omitted.
\end{proof}

\begin{lemma}
\label{lemma-equivalence-categories-open-structures}
Let $X$ be a topological space.
Let $j : U \to X$ be the inclusion of an open subset.
Let $(\mathcal{C}, F)$ be a type of algebraic structure
such that $\mathcal{C}$ has an initial object $e$.
The functor
$$
j_! : \Sh(U, \mathcal{C}) \longrightarrow \Sh(X, \mathcal{C})
$$
is fully faithful. Its essential image consists exactly
of those sheaves $\mathcal{G}$ such that
$\mathcal{G}_x = e$ for all $x \in X \setminus U$.
\end{lemma}

\begin{proof}
Omitted.
\end{proof}


\begin{lemma}
\label{lemma-equivalence-categories-open-modules}
Let $(X, \mathcal{O})$ be a ringed space.
Let $j : (U, \mathcal{O}|_U) \to (X, \mathcal{O})$
be an open subspace.
The functor
$$
j_! : \textit{Mod}(\mathcal{O}|_U) \longrightarrow \textit{Mod}(\mathcal{O})
$$
is fully faithful. Its essential image consists exactly
of those sheaves $\mathcal{G}$ such that
$\mathcal{G}_x = 0$ for all $x \in X \setminus U$.
\end{lemma}

\begin{proof}
Omitted.
\end{proof}

\begin{remark}
\label{remark-j-shriek-not-exact}
Let $j : U \to X$ be an open immersion of topological spaces as above.
Let $x \in X$, $x \not \in U$. Let $\mathcal{F}$ be a sheaf of sets
on $U$. Then $j_!\mathcal{F}_x = \emptyset$ by Lemma \ref{lemma-j-shriek}.
Hence $j_!$ does not transform a final object of $\Sh(U)$
into a final object of $\Sh(X)$ unless $U = X$.
According to our conventions in
Categories, Section \ref{categories-section-exact-functor}
this means that the functor $j_!$ is not left exact
as a functor between the categories of sheaves of sets.
It will be shown later that $j_!$ on abelian sheaves is exact,
see Modules, Lemma \ref{modules-lemma-j-shriek-exact}.
\end{remark}








